% =========================================================
% 01_vectors_foundations.tex
% Vectors and Linear Algebra
% =========================================================

% =========================================================
% SECTION DIVIDER
% =========================================================

\section{Vectors and Coordinate Systems}

\begin{frame}{Vectors and Coordinate Systems}

\begin{center}

\vspace{0.5cm}

{\Large\color{PastelBlue}\textbf{Vectors}}

\vspace{0.2cm}

{\large
Objects used to represent
magnitude, direction, position, and features
}

\vspace{0.8cm}

\begin{tikzpicture}[
    >=Latex,
    scale=0.9
]

% Axes
\draw[->, thick, gray] (-3,0) -- (3.5,0) node[right] {$x$};
\draw[->, thick, gray] (0,-2) -- (0,2.5) node[above] {$y$};

% Vector
\draw[->, very thick, PastelBlue]
    (0,0) -- (2.4,1.6)
    node[midway, above] {$\vect{v}$};

% Point
\fill[PastelLavender] (2.4,1.6) circle (2pt);
\node[above right] at (2.4,1.6) {$\vect{v}=(v_1,v_2)$};

\end{tikzpicture}

\end{center}

\end{frame}


% =========================================================
% WHAT IS A VECTOR?
% =========================================================

\begin{frame}{What Is a Vector?}

A \textbf{vector} is a mathematical object that has:

\begin{itemize}
    \item magnitude
    \item direction
\end{itemize}

\vspace{0.3cm}

A vector in two-dimensional space can be represented as

\[
\vect{v}
=
\begin{bmatrix}
v_1\\
v_2
\end{bmatrix}.
\]

\vspace{0.3cm}

\begin{block}{Key Idea}
The entries of a vector are its \textbf{components} with respect to a chosen coordinate system.
\end{block}

\end{frame}


% =========================================================
% COORDINATE SYSTEM
% =========================================================

\begin{frame}{Coordinate Systems}

A coordinate system provides a reference for describing vectors and points.

\vspace{0.2cm}

For the standard two-dimensional Cartesian system:

\begin{columns}[c,totalwidth=\textwidth]

% ---------- LEFT COLUMN ----------
\begin{column}{0.40\textwidth}

\begin{center}

\raisebox{0.60cm}{%
$\displaystyle
\vect{v}
=
\begin{bmatrix}
v_x\\
v_y
\end{bmatrix}
$
}

\end{center}

\end{column}

% ---------- RIGHT COLUMN ----------
\begin{column}{0.60\textwidth}

\begin{center}

\begin{tikzpicture}[
    >=Latex,
    scale=0.75
]

% Grid
\draw[step=1cm, gray!20, very thin]
    (-3,-2) grid (3,3);

% Axes
\draw[->, thick, gray]
    (-3,0) -- (3.5,0)
    node[right] {$x$};

\draw[->, thick, gray]
    (0,-2) -- (0,3.5)
    node[above] {$y$};

% Vector
\draw[->, very thick, PastelBlue]
    (0,0) -- (2,2.5);

% Point
\fill[PastelLavender]
    (2,2.5) circle (2.5pt);

% Vector label
\node[above right] at (2,2.5)
    {$\vect{v}=(v_x,v_y)$};

\end{tikzpicture}

\end{center}

\end{column}

\end{columns}

\end{frame}

% =========================================================
% VECTOR AS A COLUMN
% =========================================================

\begin{frame}{Representing Vectors}

Vectors are commonly written as column vectors:
\[
\vect{v}
=
\begin{bmatrix}
v_1\\
v_2\\
v_3
\end{bmatrix}
\in \mathbb{R}^3.
\]

More generally,
\[
\vect{v}
=
\begin{bmatrix}
v_1\\
v_2\\
\vdots\\
v_n
\end{bmatrix}
\in \mathbb{R}^n.
\]

\begin{block}{Interpretation}
A vector in $\mathbb{R}^n$ has $n$ components and can be viewed as a point or a directed quantity in an $n$-dimensional space.
\end{block}

\end{frame}


% =========================================================
% VECTOR ADDITION
% =========================================================

\begin{frame}{Vector Addition}

Two vectors of the same dimension can be added component-wise.

\[
\vect{u}
=
\begin{bmatrix}
u_1\\
u_2
\end{bmatrix},
\qquad
\vect{v}
=
\begin{bmatrix}
v_1\\
v_2
\end{bmatrix}
\]

give

\[
\boxed{
\vect{u}+\vect{v}
=
\begin{bmatrix}
u_1+v_1\\
u_2+v_2
\end{bmatrix}
}
\]

\vspace{0.4cm}

\begin{exampleblock}{Example}
\[
\begin{bmatrix}
2\\
1
\end{bmatrix}
+
\begin{bmatrix}
3\\
4
\end{bmatrix}
=
\begin{bmatrix}
5\\
5
\end{bmatrix}
\]
\end{exampleblock}

\end{frame}


% =========================================================
% GEOMETRIC VECTOR ADDITION
% =========================================================

\begin{frame}{Geometric Interpretation of Addition}

Vector addition can be visualized using the \textbf{parallelogram rule}.

\vspace{0.2cm}

\begin{center}

\begin{tikzpicture}[
    >=Latex,
    scale=0.85
]

\coordinate (O) at (0,0);
\coordinate (U) at (3,1);
\coordinate (V) at (1,2);
\coordinate (W) at (4,3);

\draw[->, very thick, PastelBlue]
    (O) -- (U)
    node[midway, below] {$\vect{u}$};

\draw[->, very thick, PastelLavender]
    (O) -- (V)
    node[midway, left] {$\vect{v}$};

\draw[dashed, gray]
    (U) -- (W);

\draw[dashed, gray]
    (V) -- (W);

\draw[->, very thick]
    (O) -- (W)
    node[midway, above] {$\vect{u}+\vect{v}$};

\fill (O) circle (1.5pt);
\fill (W) circle (1.5pt);

\end{tikzpicture}

\end{center}

\end{frame}


% =========================================================
% SCALAR MULTIPLICATION
% =========================================================

\begin{frame}{Scalar Multiplication}

A \textbf{scalar} is a number that can scale a vector.

For a scalar $c$,

\begin{columns}[c,totalwidth=\textwidth]

% ---------- LEFT ----------
\begin{column}{0.45\textwidth}

\begin{center}

\[
\vect{v}
=
\begin{bmatrix}
v_1\\
v_2\\
\vdots\\
v_n
\end{bmatrix}
\]

\end{center}

\end{column}

% ---------- RIGHT ----------
\begin{column}{0.45\textwidth}

\begin{center}

\[
c\vect{v}
=
\begin{bmatrix}
cv_1\\
cv_2\\
\vdots\\
cv_n
\end{bmatrix}
\]

\end{center}

\end{column}

\end{columns}

\vspace{0.4cm}

\begin{itemize}
    \item $|c|>1$   : vector is stretched
    \item $0<|c|<1$ : vector is shortened
    \item $c<0$     : direction is reversed
\end{itemize}

\end{frame}

% =========================================================
% VECTOR MAGNITUDE
% =========================================================

\begin{frame}{Magnitude of a Vector}

The magnitude or Euclidean norm measures the length of a vector.

For
\vspace{-0.5cm}

\begin{center}

\[
\vect{v}
=
\begin{bmatrix}
v_1\\
v_2
\end{bmatrix}
\qquad
\boxed{
\|\vect{v}\|
=
\sqrt{v_1^2+v_2^2}
}
\]

\end{center}

\vspace{-0.5cm}

\noindent For $\vect{v}\in\mathbb{R}^n$,

\begin{center}

\[
\boxed{
\|\vect{v}\|
=
\sqrt{
\sum_{i=1}^{n} v_i^2
}
}
\]

\end{center}

\vspace{-0.2cm}

\begin{exampleblock}{Example}

\[
\vect{v}
=
\begin{bmatrix}
3\\
4
\end{bmatrix}
\quad\Longrightarrow\quad
\|\vect{v}\|
=
\sqrt{3^2+4^2}
=
5
\]

\end{exampleblock}

\end{frame}


% =========================================================
% DISTANCE BETWEEN VECTORS
% =========================================================

\begin{frame}{Distance Between Vectors}

The distance between two vectors is the magnitude of their difference.

For $\vect{u},\vect{v}\in\mathbb{R}^n$,

\[
\boxed{
d(\vect{u},\vect{v})
=
\|\vect{u}-\vect{v}\|
}
\]

\vspace{0.5cm}

For example,

\[
\vect{u}
=
\begin{bmatrix}
3\\
4
\end{bmatrix},
\qquad
\vect{v}
=
\begin{bmatrix}
1\\
1
\end{bmatrix}
\]

give

\[
d(\vect{u},\vect{v})
=
\left\|
\begin{bmatrix}
2\\
3
\end{bmatrix}
\right\|
=
\sqrt{13}.
\]

\end{frame}


% =========================================================
% VECTORS IN COMPUTING
% =========================================================

\begin{frame}{Vectors in Computing}

Vectors provide a common representation for numerical information.

\vspace{0.3cm}

\begin{columns}[T,totalwidth=\textwidth]

\begin{column}{0.31\textwidth}

\begin{block}{Machine Learning}
Feature vectors represent measurable properties of an observation.
\end{block}

\end{column}

\begin{column}{0.31\textwidth}

\begin{block}{Computer Graphics}
Coordinates and transformations can be represented using vectors.
\end{block}

\end{column}

\begin{column}{0.31\textwidth}

\begin{block}{Data Representation}
Objects such as documents or records can be represented numerically.
\end{block}

\end{column}

\end{columns}

\vspace{0.5cm}

\begin{center}
\[
\boxed{
\text{Real-world information}
\quad\longrightarrow\quad
\text{Numerical vector}
}
\]
\end{center}

\end{frame}